\documentclass[11pt,a4paper]{article}
% Shared preamble for the qi26_25 weekly deliverable reports.
%
% Each weekly report is a short standalone document that inputs this file, so
% that formatting stays consistent across all six without duplicating the
% package list in every one. Change something here and every report follows.
%
% Usage in a weekly report:
%     \documentclass[11pt,a4paper]{article}
%     % Shared preamble for the qi26_25 weekly deliverable reports.
%
% Each weekly report is a short standalone document that inputs this file, so
% that formatting stays consistent across all six without duplicating the
% package list in every one. Change something here and every report follows.
%
% Usage in a weekly report:
%     \documentclass[11pt,a4paper]{article}
%     % Shared preamble for the qi26_25 weekly deliverable reports.
%
% Each weekly report is a short standalone document that inputs this file, so
% that formatting stays consistent across all six without duplicating the
% package list in every one. Change something here and every report follows.
%
% Usage in a weekly report:
%     \documentclass[11pt,a4paper]{article}
%     \input{weekly_preamble}
%     \weekhead{1}{Vision Foundation and Data Encoding}{Week of ...}
%     ... body ...
%     \end{document}

\usepackage[utf8]{inputenc}
\usepackage[T1]{fontenc}
\usepackage[margin=2.5cm]{geometry}
\usepackage{amsmath,amssymb}
\usepackage{booktabs}
\usepackage{graphicx}
\usepackage{enumitem}
\usepackage{listings}
\usepackage{xcolor}
\usepackage[hidelinks]{hyperref}
\usepackage{titlesec}
\usepackage{fancyhdr}

\definecolor{codebg}{RGB}{247,247,250}
\definecolor{codekw}{RGB}{20,90,160}
\definecolor{codecm}{RGB}{110,120,130}
\definecolor{accent}{RGB}{20,90,160}

\lstset{
  backgroundcolor=\color{codebg},
  basicstyle=\ttfamily\footnotesize,
  keywordstyle=\color{codekw}\bfseries,
  commentstyle=\color{codecm}\itshape,
  breaklines=true,
  frame=single,
  rulecolor=\color{gray!30},
  language=Python,
  showstringspaces=false,
  xleftmargin=6pt,
  framexleftmargin=6pt
}

\setlist{itemsep=2pt,topsep=4pt}
\titleformat{\section}{\normalfont\large\bfseries\color{accent}}{\thesection}{0.8em}{}
\titleformat{\subsection}{\normalfont\normalsize\bfseries}{\thesubsection}{0.8em}{}

\pagestyle{fancy}
\fancyhf{}
\fancyfoot[C]{\small\thepage}
\fancyfoot[R]{\scriptsize\texttt{github.com/Pushkar0997/qi26\_25}}
\renewcommand{\headrulewidth}{0pt}

% \weekhead{number}{title}{deliverable as stated in the project brief}
\newcommand{\weekhead}[3]{%
  \begin{center}
    {\scriptsize\textsf{QINTERN 2026 \textbullet\ QWORLD \textbullet\ WEEKLY DELIVERABLE}}\\[0.25cm]
    {\LARGE\bfseries Week #1: #2}\\[0.35cm]
    {\normalsize Hybrid QML and Quantum String Algorithms for Document Extraction}\\[0.25cm]
    {\small Pushkar Kumar \quad\textbullet\quad Mentor: Potluri Krishna Priyatham}\\[0.4cm]
    \fbox{\begin{minipage}{0.88\textwidth}
      \small\textbf{Deliverable as specified in the project brief:} #3
    \end{minipage}}
  \end{center}
  \vspace{0.4cm}
}

% A boxed note for caveats and honest limitations.
\newcommand{\honestnote}[1]{%
  \vspace{0.2cm}
  \noindent\fcolorbox{gray!40}{gray!8}{%
    \begin{minipage}{0.97\textwidth}\small #1\end{minipage}}
  \vspace{0.2cm}
}

%     \weekhead{1}{Vision Foundation and Data Encoding}{Week of ...}
%     ... body ...
%     \end{document}

\usepackage[utf8]{inputenc}
\usepackage[T1]{fontenc}
\usepackage[margin=2.5cm]{geometry}
\usepackage{amsmath,amssymb}
\usepackage{booktabs}
\usepackage{graphicx}
\usepackage{enumitem}
\usepackage{listings}
\usepackage{xcolor}
\usepackage[hidelinks]{hyperref}
\usepackage{titlesec}
\usepackage{fancyhdr}

\definecolor{codebg}{RGB}{247,247,250}
\definecolor{codekw}{RGB}{20,90,160}
\definecolor{codecm}{RGB}{110,120,130}
\definecolor{accent}{RGB}{20,90,160}

\lstset{
  backgroundcolor=\color{codebg},
  basicstyle=\ttfamily\footnotesize,
  keywordstyle=\color{codekw}\bfseries,
  commentstyle=\color{codecm}\itshape,
  breaklines=true,
  frame=single,
  rulecolor=\color{gray!30},
  language=Python,
  showstringspaces=false,
  xleftmargin=6pt,
  framexleftmargin=6pt
}

\setlist{itemsep=2pt,topsep=4pt}
\titleformat{\section}{\normalfont\large\bfseries\color{accent}}{\thesection}{0.8em}{}
\titleformat{\subsection}{\normalfont\normalsize\bfseries}{\thesubsection}{0.8em}{}

\pagestyle{fancy}
\fancyhf{}
\fancyfoot[C]{\small\thepage}
\fancyfoot[R]{\scriptsize\texttt{github.com/Pushkar0997/qi26\_25}}
\renewcommand{\headrulewidth}{0pt}

% \weekhead{number}{title}{deliverable as stated in the project brief}
\newcommand{\weekhead}[3]{%
  \begin{center}
    {\scriptsize\textsf{QINTERN 2026 \textbullet\ QWORLD \textbullet\ WEEKLY DELIVERABLE}}\\[0.25cm]
    {\LARGE\bfseries Week #1: #2}\\[0.35cm]
    {\normalsize Hybrid QML and Quantum String Algorithms for Document Extraction}\\[0.25cm]
    {\small Pushkar Kumar \quad\textbullet\quad Mentor: Potluri Krishna Priyatham}\\[0.4cm]
    \fbox{\begin{minipage}{0.88\textwidth}
      \small\textbf{Deliverable as specified in the project brief:} #3
    \end{minipage}}
  \end{center}
  \vspace{0.4cm}
}

% A boxed note for caveats and honest limitations.
\newcommand{\honestnote}[1]{%
  \vspace{0.2cm}
  \noindent\fcolorbox{gray!40}{gray!8}{%
    \begin{minipage}{0.97\textwidth}\small #1\end{minipage}}
  \vspace{0.2cm}
}

%     \weekhead{1}{Vision Foundation and Data Encoding}{Week of ...}
%     ... body ...
%     \end{document}

\usepackage[utf8]{inputenc}
\usepackage[T1]{fontenc}
\usepackage[margin=2.5cm]{geometry}
\usepackage{amsmath,amssymb}
\usepackage{booktabs}
\usepackage{graphicx}
\usepackage{enumitem}
\usepackage{listings}
\usepackage{xcolor}
\usepackage[hidelinks]{hyperref}
\usepackage{titlesec}
\usepackage{fancyhdr}

\definecolor{codebg}{RGB}{247,247,250}
\definecolor{codekw}{RGB}{20,90,160}
\definecolor{codecm}{RGB}{110,120,130}
\definecolor{accent}{RGB}{20,90,160}

\lstset{
  backgroundcolor=\color{codebg},
  basicstyle=\ttfamily\footnotesize,
  keywordstyle=\color{codekw}\bfseries,
  commentstyle=\color{codecm}\itshape,
  breaklines=true,
  frame=single,
  rulecolor=\color{gray!30},
  language=Python,
  showstringspaces=false,
  xleftmargin=6pt,
  framexleftmargin=6pt
}

\setlist{itemsep=2pt,topsep=4pt}
\titleformat{\section}{\normalfont\large\bfseries\color{accent}}{\thesection}{0.8em}{}
\titleformat{\subsection}{\normalfont\normalsize\bfseries}{\thesubsection}{0.8em}{}

\pagestyle{fancy}
\fancyhf{}
\fancyfoot[C]{\small\thepage}
\fancyfoot[R]{\scriptsize\texttt{github.com/Pushkar0997/qi26\_25}}
\renewcommand{\headrulewidth}{0pt}

% \weekhead{number}{title}{deliverable as stated in the project brief}
\newcommand{\weekhead}[3]{%
  \begin{center}
    {\scriptsize\textsf{QINTERN 2026 \textbullet\ QWORLD \textbullet\ WEEKLY DELIVERABLE}}\\[0.25cm]
    {\LARGE\bfseries Week #1: #2}\\[0.35cm]
    {\normalsize Hybrid QML and Quantum String Algorithms for Document Extraction}\\[0.25cm]
    {\small Pushkar Kumar \quad\textbullet\quad Mentor: Potluri Krishna Priyatham}\\[0.4cm]
    \fbox{\begin{minipage}{0.88\textwidth}
      \small\textbf{Deliverable as specified in the project brief:} #3
    \end{minipage}}
  \end{center}
  \vspace{0.4cm}
}

% A boxed note for caveats and honest limitations.
\newcommand{\honestnote}[1]{%
  \vspace{0.2cm}
  \noindent\fcolorbox{gray!40}{gray!8}{%
    \begin{minipage}{0.97\textwidth}\small #1\end{minipage}}
  \vspace{0.2cm}
}


\begin{document}
\weekhead{4}{Quantum String Operations (Processing Stage)}
{Functional Grover-based string search simulation; oracle for $O(\sqrt{N})$ search complexity.}

\section{Objective}

Implement pattern matching using Grover-based logic: develop a quantum string
matching algorithm to find structural identifiers such as document ID numbers,
and build the oracle.

\section{A Construction That Does Not Search}

The starter implementation for this track, written in July, contained the
following:

\begin{lstlisting}
match_positions = [i for i in range(len(text) - len(pattern) + 1)
                   if text[i:i+len(pattern)] == pattern]
oracle = build_marked_state_oracle(n_qubits, match_positions)
\end{lstlisting}

The comparison executes in Python, on the first line, before a single qubit is
allocated. The circuit is then constructed to flip the phase of positions it has
already been handed.

Run it and it appears to work: the measured distribution peaks at the correct
position and the amplification is textbook. It is a correct demonstration of
Grover \emph{amplification}. It is not string matching, and no complexity claim
can rest on it, because the $O(N)$ classical scan has already happened.

\honestnote{%
This construction is common in Grover string-matching demonstrations, and it is
easy to ship without noticing---the output is indistinguishable from a working
search. Identifying it is the reason this week's deliverable was rebuilt rather
than extended.}

\section{An In-Circuit Comparator Oracle}

The replacement receives only the text and the pattern. Two registers:

\begin{itemize}
\item \texttt{pos}: $\lceil\log_2(N-M+1)\rceil$ qubits, holding candidate start
      positions in superposition.
\item \texttt{win}: $M \times b$ qubits, where $b$ is the alphabet width in
      bits, holding the window of text at the candidate position.
\end{itemize}

\subsection{Four Stages}

\textbf{1. Window loading.} For each concrete position $i$, multi-controlled
writes load the bits of $\text{text}[i:i+M]$ into \texttt{win}, conditioned on
$\texttt{pos}=i$:
\begin{equation}
|i\rangle|0\rangle \;\longrightarrow\; |i\rangle|\text{text}[i:i+M]\rangle
\end{equation}
Only 1-bits require a gate, since the register starts at $|0\rangle$.

\textbf{2. Pattern XOR.} An $X$ gate is applied to $\texttt{win}[j]$ wherever
pattern bit $j$ is 1. Because the pattern is classical and known at circuit-build
time, this needs no controls. Afterwards the register holds
$\text{window} \oplus \text{pattern}$, which is all-zeros exactly when they
match---converting ``are these strings equal?'' into ``is this register
all-zeros?'', a condition a single multi-controlled gate can test.

\textbf{3. Phase marking.} Multi-controlled gates trigger on all-ones, so the
register is sandwiched between $X$ layers: flip everything, test, flip back.

\textbf{4. Uncomputation.} Stages 2 and 1 are reversed. Both are built purely
from $X$ and multi-controlled $X$ gates, each its own inverse, so the same
routine serves both directions.

\subsection{Why Uncomputation Is Mandatory}

After stage 1, \texttt{win} is \emph{entangled} with \texttt{pos}: state
$|i\rangle$ is paired with the text at position $i$. The phase from stage 3
attaches to that joint state.

Stopping there and applying the diffusion operator to \texttt{pos} alone fails:
the window acts as a which-path record, the amplitudes cannot interfere, and
Grover amplifies nothing. Uncomputation returns \texttt{win} to $|0\rangle$ for
every branch, disentangling it so the phase belongs to \texttt{pos} alone.

\subsection{An Encoding Detail}

Alphabet codes begin at 1, never 0. When $N-M+1$ is not a power of two, the
position register addresses more states than there are valid positions. Those
unreachable states never trigger a controlled load, so their window remains
all-zeros---which is exactly the pattern stage 3 marks as a match. Reserving code
0 eliminates this failure mode entirely.

\section{Staged Verification}

Correctness was established stage by stage rather than by observing a plausible
final output, because a search returning the right answer for the wrong reason
looks identical to one that works.

\begin{table}[h]
\centering
\caption{Verification results.}
\begin{tabular}{@{}rlc@{}}
\toprule
Stage & Check & Result \\
\midrule
1 & \texttt{win} holds $\text{text}[i:i+M]$ for a probe position & pass \\
2 & Exactly the true match positions have negative amplitude & pass \\
3 & Uncomputation leakage (probability outside $\texttt{win}=0$) & $\mathbf{0.00}$ \\
4 & Full search returns a true match & pass, 95--96\% \\
\bottomrule
\end{tabular}
\end{table}

On a 16-character hexadecimal field with a 2-character pattern: 14 qubits, 3
Grover iterations, correct position recovered in 95.2\% of shots.

The zero-leakage measurement is load-bearing beyond this week: it is what
licenses the browser demonstration to simulate the position register alone and
have that reduction be exact rather than approximate.

\section{Over-Rotation}

A common intuition is that more Grover iterations sharpen the result. They do the
opposite: amplitude rotates \emph{past} the target and back down. Success
probability oscillates rather than converging, and the optimal count
$\lfloor\frac{\pi}{4}\sqrt{N/M}\rfloor$ is a genuine optimum, not a minimum.

\section{Resource Scaling}

Grover provides $O(\sqrt{N})$ \emph{queries}. But the window-loading stage must
write the text into the circuit, and each candidate position requires its own
controlled load.

\begin{table}[h]
\centering
\caption{Measured oracle resources, hexadecimal alphabet, 2-character pattern.}
\begin{tabular}{@{}rrrr@{}}
\toprule
$N$ (characters) & Qubits & Depth & CX gates \\
\midrule
8 & 13 & 1{,}857 & 1{,}082 \\
16 & 14 & 6{,}383 & 3{,}428 \\
32 & 15 & 16{,}159 & 8{,}570 \\
64 & 16 & 39{,}558 & 20{,}904 \\
\bottomrule
\end{tabular}
\end{table}

Fitting a power law in log--log space gives a CX scaling exponent of
$\mathbf{1.39}$---\emph{superlinear} in text length. The measured curve lies
above a linear reference and nowhere near $\sqrt{N}$.

\honestnote{%
\textbf{The $O(\sqrt{N})$ advantage does not survive.} For searching stored,
unstructured classical text that must be loaded into the circuit, the loading
cost dominates the search saving. This is a known caveat in the literature; the
contribution here is reproducing it by direct measurement rather than quoting
$\sqrt{N}$ unqualified.

Grover's advantage remains real where the text is already available as a quantum
oracle or in QRAM; where the search space is \emph{generated} by a function
rather than stored, as in satisfiability or cryptographic preimage search; or
where a loaded state is reused across many queries. None of these describes
scanning a document just received.}

\section{A Second Practical Limit}

A full 38-symbol alphabet requires 6 qubits per character, so even a short search
window exceeds what is simulable on ordinary hardware. The pipeline therefore
searches the hexadecimal-representable identifier field, which is the realistic
target for structured-field lookup in any case.

\begin{table}[h]
\centering
\caption{Alphabet width against window cost.}
\begin{tabular}{@{}lrr@{}}
\toprule
Alphabet & Bits per character & Qubits for an 8-character window \\
\midrule
Binary (2 symbols) & 1 & 8 \\
Hexadecimal (16) & 5 & 40 \\
Full charset (38) & 6 & 48 \\
\bottomrule
\end{tabular}
\end{table}

\section{Deliverable Status}

Complete. A functional Grover-based string search that genuinely searches,
verified stage by stage with zero uncomputation leakage, with resource scaling
measured rather than assumed.

\vspace{0.3cm}
\noindent\textbf{Artifacts:} \texttt{track\_b\_search/oracle/comparator\_oracle.py},
\texttt{notebooks/03\_grover\_comparator.ipynb}

\vspace{0.2cm}
\noindent\textbf{Self-test:} \texttt{python track\_b\_search/oracle/comparator\_oracle.py}

\end{document}
